\scanpage{14}
\mspage{II -18 -}\hfill n$^\circ$338

COROLLAIRE. - \uline{Soient} $m,m'\handwritten{\in}M$, \uline{et} $s,s'\handwritten{\in}\bar S$. \uline{Pour que} $m/s=m'/s'$ \uline{dans} $M[S^{-1}]$, \uline{il faut et il suffit qu'il existe} $t\handwritten{\in}\bar S$ \uline{vérifiant} $t(sm'-s'm)=0$.

En effet les relations $m/s=m'/s'$, $(sm'-s'm)/ss'=0$ et $(sm'-s'm)/1=0$ sont équivalentes ; il suffit alors d'appliquer la prop. 3.

PROPOSITION 4. - \uline{Soient} $A$ \uline{et} $B$ \uline{deux anneaux}, $S$ \uline{une partie de} $A$, $T$ \uline{une partie de} $B$, \uline{et} $f$ \uline{un homomorphisme de} $A$ \uline{dans} $B$ \uline{tels que} $f(S)\handwritten{\subset}T$. \uline{Soient} $M$ \uline{un} $A$-\uline{module}, $\unsure{N|N'}$ \uline{un} $B$-\uline{module}, \uline{et} $u$ \uline{une application} $A$-\uline{linéaire de} $M$ \uline{dans} $N$. \uline{Il existe alors une application} $(A[S^{-1}])$-\uline{linéaire} $u'$ \uline{et une seule de} $M[S^{-1}]$ \uline{dans} $N[T^{-1}]$ \uline{telle que} $u'(m/1)=u(m)/1$ \uline{pour tout} $m\handwritten{\in}M$.

En effet l'application $i_N^T\handwritten{\circ}u$ de $M$ dans $N[T^{-1}]$ est $A$-linéaire. En outre, si $s\handwritten{\in}S$, on a $f(s)\handwritten{\in}T$, donc $s$ définit dans $N[T^{-1}]$ une homothétie bijective. L'existence et l'unicité de $u'$ résultent alors de la définition de $M[S^{-1}]$ comme solution de problème d'application universelle.

Avec les mêmes notations, soient $C$ un troisième anneau, $U$ une partie de $C$, $g$ un homomorphisme de $B$ dans $C$ tels que $g(T)\handwritten{\subset}U$, $P$ un $C$-module, $v$ une application $B$-linéaire de $N$ dans $P$, et $v'$ l'application $(B[T^{-1}])$-linéaire de $N[T^{-1}]$ dans $P[U^{-1}]$ associée à $v$. On a alors
\[
(v\handwritten{\circ}u)'=v'\handwritten{\circ}u'\,.\tag*{(3)}
\]
De même, avec des notations évidentes, on a
\[
(u+u_1)'=u'+u'_1\,.\tag*{(4)}
\]

\uline{Remarques}. - 1) Si, dans la proposition 4, on prend $B=A$, $T=S$, et $f$ égal à l'application identique de $A$, on voit aussitôt que $u'$ n'est autre que l'application $1\handwritten{\otimes}u$ de $A[S^{-1}]\handwritten{\otimes}M$ dans $A[S^{-1}]\handwritten{\otimes}N$. Nous la noterons désormais $u^S$.

2) Soient $A$ un anneau, $S$ une partie de $A$, et $M$ un $A$-module. Si $T$ est une partie de $A$ contenant $S$, la proposition 4 prouve qu'il existe une application $(A[S^{-1}])$-linéaire $i_M^{T,S}$ et une seule de $M[S^{-1}]$ dans $M[T^{-1}]$ telle que
% archive sheet number: 13

